\summarychapter{37}{A New Social Frontier}{Nineteenth-Century England}

\begin{inbrief}
Nineteenth-century England changed at a pace that left many fearful and others hopeful. Two small
movements helped millions of believers adjust, and they answered the social crisis in opposite ways.
The Evangelicals of the Clapham Sect went out to reform society and ended the British slave trade
and slavery; the high-church men of the Oxford movement drew in to defend the sanctity of the
church. Between them they stand for mission and worship.
\question{How are Christians to view the world?}
\end{inbrief}

``Lead, kindly Light, amid th' encircling gloom'' was written in 1833 by John Henry Newman on his
way home from Sicily. Its somber mood belongs to many troubled souls of the time; a decade later
Newman fled to Rome for safety. The same gloom sounds in the Evangelical Henry Francis Lyte's ``Abide
with Me'': ``Change and decay in all around I see.'' Neither movement of this chapter was large at
first, a reminder of Gilbert Murray's remark that ``the uplifting of man has been the work of a
chosen few.'' Evangelicals still look to Clapham as a model of social concern, and high-church
Anglicans to Oxford as a wellspring of devout churchmanship.

\ssection{Evangelicals in the World}
The church has a double commission. God sends his people into the world to proclaim salvation and
serve the needy, and he calls them out of the world, to adore him and be taught by him. Without worship, mission may turn into hollow service, and worship without mission careless religion, so the church is
always saying yes here and no there. In England, where society changed so fast, Protestants could not always tell friend from foe among those they addressed.

\begin{sidebox}{Evangelicals}
Current descriptions often rest on David Bebbington's four characteristics: conversionism, activism,
biblicism (confidence in the Bible as the source of devotion, guidance and theology) and
crucicentrism (stress on the sacrificial character of the cross). Thomas Kidd adds an emphasis on the
Holy Spirit, shown in revivals, outpourings of the Spirit and meetings with God's love that change a life.
\end{sidebox}

The century was in many respects a British one. The Industrial Revolution had begun in England, and London was the largest city on earth and its financial center, British commerce and the navy were
everywhere, and by 1914 Britain ruled the largest empire ever made. The growth left many breathless.
Some, remembering the French Revolution, feared the future; others hailed change as progress. English
Protestants were either in the established Anglican church or in the Nonconformist bodies:
Methodists, Baptists, Congregationalists and a few smaller ones. But the striking movements did not
run along denominational lines. The liberties of the age let Christians form a host of religious
societies for particular objects, such as distributing Bibles or relieving the poor, groups of
individuals with no sacraments, creeds or ordained ministers of their own.

The greatest power in English religion was the evangelical movement that came from Wesley and
Whitefield. Its marks were intense personal piety, usually springing from a conversion, and energetic
service in the world, both fed by devotion to the Bible and directed by the themes of the revival:
God's love in Christ, salvation through faith, and the new birth by the Holy Spirit. A significant
minority of Anglican pulpits and a majority of Nonconformist ones preached it. Anglican Evangelicals
were loyal to their church and its bishops but glad to work with Nonconformists, because they thought
the preaching of the gospel more important than sacraments or styles of ritual, a position called
``Low Church.'' They saw the ills of British society as a call to service.

\ssection{The Clapham Community}
Their headquarters was Clapham, then a village three miles from London and the country home of a
group of wealthy and ardent Evangelicals who practiced ``saintliness in daily life.'' Historians call
them the Clapham Sect, but they were less a sect than a close family. The parish minister, John Venn,
was their spiritual guide, and they met for Bible study, talk and prayer in the oval library of the
banker Henry Thornton. Their unquestioned leader was the parliamentarian William Wilberforce
(1759--1833). Around him were John Shore (Lord Teignmouth), who had governed India; Charles Grant, who chaired the East India Company; James Stephen, undersecretary for the colonies; Zachary Macaulay,
editor of the \emph{Christian Observer}; the abolitionist Thomas Clarkson; and others.

Wilberforce had been converted at twenty-five through a book, \emph{Rise and Progress of Religion in the Soul} by Philip Doddridge. He had wealth, a liberal education and unusual gifts. William Pitt
said he had the greatest natural eloquence he had known, and some called him ``the nightingale of the
House of Commons''; many spoke of his gift for friendship and his principles. ``My business is in the
world,'' he said, ``and I must mix in the assemblies of men or quit the part which Providence seems
to have assigned me.'' In what they called their ``Cabinet Councils'' the friends discussed the
country's wrongs and the battles to be fought, and then in and out of Parliament acted as one body,
each doing what he did best. His biographer Reginald Coupland says there has been nothing like it
since in British public life.

\ssection{Evangelicals and Social Issues}
From Clapham came a stream of societies: for missions (the Church Missionary Society, 1799), for Bibles (the British and Foreign Bible Society, 1804), for relief of the poor and for prison reform. Its greatest effort was the campaign against slavery, beginning with the
slave trade, the capture of Africans for sale in the West Indies. England had entered the trade in
1562, when Sir John Hawkins sold a cargo from Sierra Leone in St.~Domingo. After 1660 Charles II
chartered a company that carried three thousand slaves a year, and by 1770 British ships carried more
than half of the hundred thousand taken yearly from West Africa. Many thought the trade inseparable
from the nation's commerce and security.

Wilberforce first spoke against it in the Commons in 1789. He saw that eloquence alone would not
defeat the commercial interests and turned to his friends for reliable facts. Two years later, after
exhaustive preparation, he sought leave for a bill to stop the import of slaves into the West Indies:
``Never, never will we desist till we have wiped away this scandal from the Christian name.'' Oratory
again fell short, but support grew, and the abolitionists saw that they must appeal to the people as
well as to Parliament. Step by step they learned the two things democratic politics requires: shaping public opinion and then turning it into pressure on those who govern. They gathered petitions, published
good literature, lectured, and campaigned on billboards. Nonconformists rallied to them, and for the
first time women took part in a political contest. Then Wilberforce and four Clapham colleagues, the
``Saints'' in the Commons, carried the fire into Parliament.

\ssection{The End of Slave Trade}
On February 23, 1807, the opposition broke. When one member contrasted Wilberforce with Napoleon,
the House rose and cheered as it seldom had, while Wilberforce sat bowed with his head in his hands
and tears on his face. The legal trade was ended, but the slaves were still in chains. He fought on
for emancipation until age and ill health took him from Parliament, and then enlisted a young
Evangelical, Thomas Fowell Buxton, to lead the ``holy enterprise.'' Passage of the Emancipation Act,
freeing the slaves throughout the British Empire, was assured on July 25, 1833, four days before
Wilberforce died. Coming before the European powers partitioned Africa, the act mattered enormously.
G.~M.~Trevelyan wrote that on the last night of slavery the West Indian slaves climbed the hills to
watch the sun rise on their freedom, while in the forests of Africa no one marked the day; yet it
was Africa that was most deeply affected, for the most powerful of the nations that would shape its future had resolved that black and white should not meet as slave and master. It is above all for this that people still point to the Clapham Sect when they want to show how a few able and devoted people can influence a
society, perhaps the world.

\ssection{The Oxford Movement}
The Oxford movement was a contrasting response. Like the evangelical movement it was more of the
heart than of the head, but its leaders watched the course of English society with alarm. To them the government's reforms were attacks on the sanctity of the Church of England and resolved to
resist the world's intrusion. ``Now each man attempts to judge for himself,'' Newman wrote to his
mother in 1829; ``the talent of the day is against the Church.'' The church's strength had rested on
landed aristocrats powerful in Parliament, while the new industrial towns such as Manchester and
Birmingham had no representatives. The Reform Act of 1832 shifted power from the gentry to the middle
class, so that many members of Parliament who were not Anglicans now had power over the church. Were
profane politicians to lay hands on holy things?

On July 14, 1833, John Keble, fellow of Oriel College, preached at Oxford on ``National Apostasy'': a
nation denies God's sovereignty when it shows disrespect to the successors of the apostles, the
bishops, and appeals only to popularity or expediency. He was joined by John Henry Newman
(1801--1890), the University Church's vicar, and by Edward Pusey, professor of Hebrew, and the three
turned a protest into a movement. The church's authority, they held, does not come from the state but
from God, and bishops hold office by apostolic commission, not social position; even wholly separated
from the state, the Church of England could claim Englishmen's allegiance. In 1833 they began the
``Tracts for the Times,'' whence the name Tractarians. The tracts expounded one article of the Creed,
belief in ``one, holy, Catholic and apostolic Church.'' They stressed the apostolic succession of
bishops and the church's God-given authority to teach truth and rule lives, ascribed real saving
power to the sacraments, and held up as the ideal the undivided church of the first five centuries.
Some of this history was fanciful, but they believed it ardently. They called themselves Catholics
and shunned the name Protestant because it spoke of division. Public worship was vital to them.
Symbolic acts such as turning to the altar, kneeling and raising the cross have religious value, and
because worship asks for the response of the whole person it should appeal to the senses, with rich
vestments, incense and trained voices. Tractarianism was High Church religion at its most zealous, and the Oxford men drew closer to Rome with each step. Then came the thunderclap. In 1841 Newman's Tract 90
argued that the Thirty-nine Articles are not necessarily Protestant and can be read in a Catholic
sense. In the storm that followed, the bishop of Oxford forbade him to publish any more. Newman decided that no one could be fully Catholic outside the Church of Rome, which he joined in 1845, and in the next six years hundreds of Anglican clergy followed. Dublin's new Catholic University made him its rector, and in 1877 he was named a cardinal. Most Tractarians, however, stayed, and more and
more clergy adopted their views. For many, religion centered on ritual, priesthood and sacraments,
and the concern for beauty improved church architecture, music and art. In time people spoke of Anglo-Catholics rather than of Tractarians or the Oxford movement, meaning Anglicans who value their unity with the catholic
tradition of Eastern Orthodoxy and Rome but accept the supremacy of neither patriarch nor pope.

Both views of Christianity's place in society are still alive, if not always well. Few generations
have a Wilberforce or a Newman, but their convictions last because mission and worship are basic to
Christianity in any age. Jesus prayed for his disciples not that they be taken out of the world but
that they be kept from the evil one: ``As you sent me into the world, I have sent them into the
world.''

\begin{keyfacts}
\keypeople{William Wilberforce; John Venn; Henry Thornton; Zachary Macaulay; Thomas Clarkson; Thomas
Fowell Buxton; John Keble; John Henry Newman; Edward Pusey; Henry Francis Lyte.}
\keydates{%
  \kd{1789}{Wilberforce's first speech against the slave trade}
  \kd{1799}{Church Missionary Society founded}
  \kd{1804}{British and Foreign Bible Society founded}
  \kd{1807}{Parliament abolishes the British slave trade}
  \kd{1832}{British Reform Act shifts power to the middle class}
  \kd{1833}{Emancipation of slaves in the British Empire; death of Wilberforce; Keble's sermon and
    the first ``Tracts for the Times''}
  \kd{1841}{Newman's Tract 90}
  \kd{1845}{Newman enters the Roman Catholic Church}}
\keyterms{Clapham Sect; Evangelicals; Low Church; religious societies; abolition; emancipation; public
opinion; Oxford movement; Tractarians; apostolic succession; High Church; Anglo-Catholic; mission and
worship.}
\end{keyfacts}
